% ── 1. Project Topic ──────────────────────────────────────────────────────────
\section{Project Topic}
 
This project is about \textbf{detecting cars in snowy environments} using deep learning.
Detecting vehicles in winter conditions is harder than normal conditions because snow
can hide parts of a car, reduce contrast, and make roads and cars look similar in colour.
This makes it an interesting and challenging problem to solve.
 
We will use the \textbf{Nordic Vehicle Dataset (NVD)}, a dataset created specifically for
this purpose by researchers at LTU (Luleå University of Technology). The dataset contains
videos and images of vehicles captured in cold, snowy Nordic weather conditions.
More information about the dataset can be found at:
\url{https://nvd.ltu-ai.dev/}
 
% ── 2. Intended Approach ─────────────────────────────────────────────────────
\section{Intended Approach}
 
\subsection{Type of Problem}
 
This is an \textbf{object detection} problem. Object detection means we want the model to:
\begin{enumerate}
    \item Find \textit{where} a car is in an image (draw a bounding box around it).
    \item Confirm \textit{what} it is (classify it as a ``car'' or ``vehicle'').
\end{enumerate}
Unlike simple image classification (which only says ``is there a car?''), object detection
tells us both the location and the identity of each object in the image.
 
\subsection{Dataset Split}
 
We will follow the official data split provided in the project description:
 
\begin{center}
\begin{tabular}{>{\bfseries}ll}
\toprule
Split & Recording \\
\midrule
Train      & 2022-12-04 Hjenberg 02 \\
           & 2022-12-23 Asjo 01\_HD 5x stab \\
           & 2022-12-02 Asjo 01\_stabilized \\
Validation & 2022-12-03 Nyland 01\_stabilized \\
Test       & 2022-12-23 Bjenberg-02 \\
\bottomrule
\end{tabular}
\end{center}
 
\subsection{Model Architecture}
 
We will use \textbf{YOLOv9} (You Only Look Once, version 9) as our base model.
YOLO is a popular family of object detection models that is:
\begin{itemize}
    \item \textbf{Fast} -- it processes the whole image in one pass (hence ``You Only Look Once'').
    \item \textbf{Accurate} -- recent versions like YOLOv9 achieve very good detection results.
    \item \textbf{Well-supported} -- there are pre-trained versions we can fine-tune on our data.
\end{itemize}
 
We plan to start with a \textbf{pre-trained YOLOv9 model} (trained on a large general dataset
such as COCO) and then \textbf{fine-tune} it on the NVD training split. Fine-tuning means
we take a model that already knows what cars look like in normal conditions and teach it
to also recognise them in snowy conditions. This usually works better and faster than
training from scratch.
 
\subsection{Loss Function}
 
YOLOv9 uses a combination of loss functions to train the model. In simple terms:
 
\begin{itemize}
    \item \textbf{Bounding box regression loss} -- penalises the model when the predicted box
          does not match the true box. We will use \textit{CIoU loss} (Complete Intersection
          over Union), which measures how well two boxes overlap.
    \item \textbf{Classification loss} -- penalises the model when it misidentifies an object.
          We will use \textit{Binary Cross-Entropy (BCE) loss}.
    \item \textbf{Objectness loss} -- penalises the model when it misses an object or detects
          something that is not there. This is also BCE loss.
\end{itemize}
 
These losses are added together during training to guide the model to make better predictions.
 
\subsection{Evaluation Metrics}
 
We will measure how well our model works using the following standard metrics:
 
\begin{itemize}
    \item \textbf{mAP (mean Average Precision)} -- the main metric for object detection.
          It summarises how precise the model is across different detection thresholds.
          We will report \textbf{mAP@0.5} (IoU threshold of 0.5) and
          \textbf{mAP@0.5:0.95} (averaged over multiple thresholds).
 
    \item \textbf{Precision} -- of all boxes the model predicts, what fraction are correct?
 
    \item \textbf{Recall} -- of all real cars in the images, what fraction did the model find?
 
    \item \textbf{F1-Score} -- the harmonic mean of Precision and Recall, giving a single
          balanced score.
 
    \item \textbf{Inference speed (FPS)} -- how many frames per second the model can process,
          which matters for real-time applications like dashcams or traffic cameras.
\end{itemize}
 
\subsection{Implementation Plan}
 
\begin{enumerate}
    \item \textbf{Week 1} -- Download and explore the NVD dataset. Understand the annotation
          format and visualise some examples with bounding boxes.
    \item \textbf{Week 2} -- Set up the YOLOv9 environment (PyTorch), prepare the dataset in
          YOLO format, and run a baseline experiment.
    \item \textbf{Week 3} -- Fine-tune the model, experiment with data augmentation (e.g.\
          adding artificial snow, brightness adjustments) to improve robustness.
    \item \textbf{Week 4} -- Evaluate on the test set, analyse failure cases, and prepare
          the final presentation and GitHub repository.
\end{enumerate}
 
\subsection{Tools and Frameworks}
 
We plan to use:
\begin{itemize}
    \item \textbf{Python} with \textbf{PyTorch} as the main deep learning framework.
    \item \textbf{YOLOv9} official implementation from GitHub.
    \item \textbf{Roboflow / custom scripts} for dataset preparation and annotation conversion.
    \item \textbf{GitHub} for version control and sharing the final code, model, and results.
\end{itemize}
 
% ── End ───────────────────────────────────────────────────────────────────────
\vspace{2em}
\hrule
\vspace{0.5em}
\footnotesize
\textit{Reference: Mokayed et al., ``Nordic Vehicle Dataset (NVD): Performance of Vehicle Detectors
Using Newly Developed Dataset,'' CVPRW 2023.
\url{https://openaccess.thecvf.com/content/CVPR2023W/AICity/html/Mokayed_Nordic_Vehicle_Dataset_NVD_Performance_of_Vehicle_Detectors_Using_Newly_CVPRW_2023_paper.html}}
 